\section{Discussion}
The results support the central hypothesis of this paper: the randomized-exploration behavior of Thompson sampling does not need to be abandoned under heavy-tailed rewards, provided the likelihood update is made robust to extreme observations. Two design choices appear important. First, truncating rather than clipping-and-discarding preserves a usable gradient signal even from large-magnitude rewards, which matters when informative rewards and outliers are not easily distinguishable ex ante. Second, calibrating the truncation threshold online, rather than fixing it from a conservative worst-case bound on $u$, avoids the over-truncation that would otherwise slow learning in the early rounds before the moment estimate stabilizes.

A limitation of our analysis is that the regret bound in Equation~\ref{eq:regret-bound} assumes $\epsilon$ is known; in practice we estimate it jointly with $u$, and our experiments suggest the resulting bound is empirically loose by roughly a constant factor rather than a different rate, but a fully data-driven analysis of joint $(\epsilon, u)$ estimation remains open. We also assume linear reward structure; extending the calibrated-truncation idea to nonlinear or kernelized reward models, in the spirit of \citet{pendelton2022sparse}, is a natural next step. Finally, all experiments here use synthetic noise; validating the calibration mechanism on reward logs with empirically heavy-tailed behavior (e.g.\ ad-auction revenue) is left to future work.
